\documentclass{report}
\usepackage{amsmath,amssymb}
\setlength{\parindent}{0mm}
\setlength{\parskip}{1em}
\begin{document}
\begin{center}
\rule{6in}{1pt} \
{\large Julianne Chung \\
{\bf Optimal Regularized Inverse Matrices as Preconditioners for Inverse Problems}}

225 Stanger Street \\ McBryde 460 \\ Blacksburg \\ VA 24061
\\
{\tt jmchung@vt.edu}\\
Matthias Chung\end{center}

Iterative methods are often used to solve inverse problems, but many
iterations may be needed to obtain a reasonable solution due to slow
convergence of the method. Preconditioners have been proposed for
accelerating the convergence of iterative methods, but inverse problems,
in particular ill-posed inverse problems, are notoriously difficult to
precondition. Optimal regularized inverse matrices were recently proposed
and studied for computing solutions to inverse problems, but in this
talk, we investigate the use of optimal regularized inverse matrices as
preconditioners for solving inverse problems. We consider an alternating
optimization update approach for Bayes risk minimization and show that
our approach can not only accelerate convergence of iterative methods,
but also stabilize the inversion since it inherently incorporates
regularization. Some theoretical results provide insight into the
approach, and numerical results on image processing examples demonstrate
the benefits of the preconditioner.


\end{document}
